\chapter{Valproic Acid Level}
\section*{What Is This Test?}
The valproic acid (VPA) level measures the serum concentration of valproic acid (also known as valproate or divalproex sodium), a branched-chain carboxylic acid that is a broad-spectrum anticonvulsant and mood stabilizer. VPA is used for: generalized absence (petit mal) seizures (drug of choice), myoclonic seizures, generalized tonic-clonic seizures, focal seizures, bipolar I disorder (acute mania and maintenance), and migraine prophylaxis. VPA has multiple mechanisms of action: potentiation of GABA-ergic neurotransmission (increasing GABA synthesis by activating glutamic acid decarboxylase and inhibiting GABA transaminase), blockade of voltage-gated sodium channels, and inhibition of histone deacetylase (HDAC), which may contribute to its mood-stabilizing and anticancer effects. VPA has a therapeutic range of 50--125 mcg/mL (total). VPA is >90\% protein-bound at therapeutic concentrations; however, protein binding is saturable, and at higher concentrations, the free fraction increases disproportionately. VPA exhibits nonlinear protein binding, meaning free VPA levels increase more than proportionally with dose increases. The free (unbound) concentration is the pharmacologically active form and may be more useful for TDM. VPA is metabolized primarily by the liver through glucuronidation (UGT enzymes) and mitochondrial beta-oxidation. Toxicity (levels >125 mcg/mL): neurotoxicity (tremor, sedation, ataxia, confusion), hepatotoxicity (idiosyncratic fulminant hepatic failure most common in children under 2 years on polytherapy, particularly with metabolic disorders --- mitochondrial hepatopathy prominent), pancreatitis (can be fatal), hyperammonemia with or without hepatic dysfunction (VPA inhibits the urea cycle enzyme carbamoyl phosphate synthetase I), thrombocytopenia (dose-dependent), teratogenicity (neural tube defects/spina bifida, 1--2\% risk with first-trimester exposure --- the highest teratogenic risk among anticonvulsants).
\section*{Alternative Names}
Valproate Level; Depakote Level; Depakene Level; VPA Level; Divalproex Sodium Level; Valproic Acid Serum Concentration; Free Valproic Acid
\section*{Principle}
Valproic acid is measured by particle-enhanced turbidimetric inhibition immunoassay (PETINIA) or enzyme-multiplied immunoassay technique (EMIT). Free VPA is measured after ultrafiltration. LC-MS/MS is the reference method with no cross-reactivity.
\section*{History}
Valproic acid was first synthesized in 1882 by B. S. Burton while investigating fatty acid derivatives and was used for decades as an organic solvent. Its anticonvulsant activity was discovered serendipitously in 1963 by Pierre Eymard and colleagues at the French laboratory Laboratoires Berthier, who observed that valproate used as a solvent for candidate anticonvulsant compounds had anticonvulsant activity of its own. Valproate was approved in France in 1967 and by the US FDA in 1978 for absence seizures. Divalproex sodium (Depakote), an enteric-coated compound with fewer gastrointestinal side effects, was approved in 1983. Its mood-stabilizing efficacy was established by the Depakote Mania Study Group trial published in 1994, leading to FDA approval for acute mania in 1995 and for migraine prophylaxis in 1996.
\section*{Why This Test Is Ordered}
TDM every 3--6 months during stable therapy, and more frequently during dose changes, suspected toxicity, pregnancy, polytherapy with enzyme-inducing drugs, and evaluating non-adherence.
\section*{Is This a Primary or Secondary Test}
TDM is a secondary monitoring test. Valproate therapy is primarily managed on the basis of seizure control, mood stability, and tolerability; the serum level is a secondary tool used to confirm therapeutic dosing, to titrate toward an effective concentration, and to evaluate suspected toxicity, non-adherence, or drug interactions. Periodic monitoring every 3--6 months during stable therapy serves as a routine safety and efficacy check rather than a diagnostic test.
\section*{How This Test Is Performed}
A blood sample is collected by standard venipuncture into a serum separator tube (SST). Total valproic acid is measured in serum by immunoassay (PETINIA or EMIT); free (unbound) valproic acid is measured after ultrafiltration of serum and reported separately. LC-MS/MS is the reference method without cross-reactivity. Results are typically available within 1--2 hours for urgent requests and 1--2 days for routine samples.
\section*{What Preparation Is Needed}
The sample should be drawn as a trough, immediately before the next scheduled dose, at steady state (2--4 days after initiation or dose change). The dose, formulation (valproic acid vs. divalproex sodium, immediate-release vs. extended-release), and time of the last dose should be documented, and the dosing schedule should be consistent. Because valproate protein binding is saturable, free (unbound) levels should be requested together with total levels in pregnancy, hypoalbuminemia, and renal or hepatic disease.
\section*{Contraindications and Precautions}
\begin{itemize}
  \item No absolute contraindications to standard venipuncture
  \item Interpret total levels with knowledge of protein binding: in pregnancy, hypoalbuminemia, and renal or hepatic disease the free fraction is increased, and total levels may underestimate the pharmacologically active drug
  \item Hepatic injury, pancreatitis, hyperammonemia, and thrombocytopenia can occur at or near therapeutic levels; liver function tests, ammonia, amylase/lipase, and platelet count are monitored along with the level
  \item Enzyme inducers (phenytoin, carbamazepine, phenobarbital) lower total valproate but can raise free levels; highly protein-bound drugs such as aspirin displace valproate and raise free levels
  \item Valproate is teratogenic (neural tube defects), so its use is managed carefully in women of childbearing potential
\end{itemize}
\section*{Risks}
Minimal risk. Slight pain or bruising at the venipuncture site. Rarely: hematoma, infection, or vasovagal syncope.
\section*{What Happens After the Test}
The result is interpreted against the therapeutic range (total 50--125 mcg/mL; free 5--25 mcg/mL) together with seizure control, mood stability, and adverse effects. Subtherapeutic levels prompt a dose increase with repeat sampling at a new steady state; toxic levels or dose-related adverse effects prompt dose reduction, rechecking of the level, and monitoring of liver function, platelets, and ammonia as indicated. A follow-up level is generally obtained 2--4 days after any dose change.
\section*{Understanding Results}
Subtherapeutic (<50 mcg/mL): Inadequate seizure control or mood stabilization. Toxic (>125 mcg/mL): Tremor (common), sedation, ataxia, confusion, hyperammonemia, thrombocytopenia. Severe toxicity (>200 mcg/mL): coma, cerebral edema, cardiovascular collapse, severe hepatotoxicity.
\section*{Normal Results}
Therapeutic range (total VPA): 50--125 mcg/mL. Free VPA therapeutic range: 5--25 mcg/mL. Trough: immediately before the next dose. Steady state: 2--4 days after initiation or dose change.
\section*{Follow-up Tests}
Liver function tests (AST, ALT, PT/INR at baseline and every 3--6 months), complete blood count with platelets (thrombocytopenia is dose-dependent, monitor every 3--6 months), serum ammonia (if altered mental status, even with normal LFTs), serum amylase/lipase (if abdominal pain/suspected pancreatitis), and serum carnitine (VPA depletes carnitine; supplementation may reduce hyperammonemia).
\section*{References}
\begin{enumerate}
  \item Loscher W. Basic pharmacology of valproate: a review after 35 years of clinical use. CNS Drugs. 2002;16(10):669--694.
  \item Perucca E. Pharmacological and therapeutic properties of valproate. CNS Drugs. 2002;16(10):695--714.
  \item Bryant AE 3rd, Dreifuss FE. Valproic acid hepatic fatalities. III. U.S. experience since 1986. Neurology. 1996;46(2):465--469.
  \item Patsalos PN, Berry DJ, Bourgeois BFD, et al. Antiepileptic drugs--best practice guidelines for therapeutic drug monitoring: a position paper by the subcommission on therapeutic drug monitoring, ILAE Commission on Therapeutic Strategies. Epilepsia. 2008;49(7):1239--1276.
  \item Bowden CL, Brugger AM, Swann AC, et al. Efficacy of divalproex vs. lithium and placebo in the treatment of mania. The Depakote Mania Study Group. JAMA. 1994;271(12):918--924.
  \item Burtis CA, Ashwood ER, Bruns DE. Tietz Textbook of Clinical Chemistry and Molecular Diagnostics. 5th ed. St. Louis, MO: Elsevier Saunders; 2012.
\end{enumerate}
\clearpage
